\documentclass[12pt]{ximera}
\typeout{Start loading xmPreamble.tex}%

% Add here extra macro's that are loaded automatically by all documents of claas 'ximera' or 'xourse' in this repo

%%
%%  Example:
%%
% \newcommand{\R}{\mathbb{R}}
\usepackage{tikz}
\usetikzlibrary{positioning, arrows.meta}
\usepackage{multicol}
\usepackage{enumitem}
\usepackage{caption}
\usepackage{amsmath}
\usepackage{amsthm}

\title{Markov Projections}
\begin{document}
\begin{abstract}
An introduction to Markov projections (lumped Markov chains), including necessary and sufficient conditions for a projection to yield a Markov chain, transition probabilities of the projected process, and the Pitman-Rogers intertwining relation.
\end{abstract}
\maketitle

\section*{Markov Projections}

\subsection*{Setup}

Assume $\{X_n : n \geq 0\}$ is a discrete-time Markov chain on a discrete state space $\mathcal{X}$, and let $\{A_y : y \in \mathcal{Y}\}$ be a finite partition of $\mathcal{X}$. We investigate when $X_n$ defines a Markov chain on this partition.

One motivation: if $\mathcal{X} = \mathbb{Z}^2$, we may only care about the $y$-coordinate of $X(t)$. For example, if $X(t)$ is a simple random walk on $\mathbb{Z}^2$ with
\[
P_{x',x} = \frac{1}{4} \cdot \mathbf{1}_{|x - x'| = 1},
\]
then the $y$-coordinate is a Markov chain on $\mathbb{Z}$ with transition probabilities
\[
P_{y',y} = \begin{cases} 1/2 & \text{if } y - y' = 0 \\ 1/4 & \text{if } |y - y'| = 1 \\ 0 & \text{if } |y - y'| > 1. \end{cases}
\]

This concept was first studied by Kemeny and Snell in \textit{Finite Markov Chains} (1960), where these were called ``lumped Markov chains.'' More generally, given a partition of $\mathcal{X}$ indexed by $\mathcal{Y}$, there is a natural projection map $\pi : \mathcal{X} \to \mathcal{Y}$.

In general, $\{Y(n) = \pi(X(n)) : n \geq 0\}$ need not be a Markov chain, because $P(Y(n+1) = y \mid Y(n) = y')$ may depend on $X(n)$. For example, if $X(t)$ is a random walk on $\mathbb{Z}^2$ that slows down farther from the $x$-axis, with $d(x)$ denoting the distance from $x$ to the $x$-axis and
\[
P_{x',x} = \begin{cases} \frac{d(x)}{d(x)+1} & \text{if } |x' - x| = 0 \\ \frac{1}{4} \cdot \frac{1}{d(x)+1} & \text{if } |x' - x| = 1 \\ 0 & \text{if } |x' - x| > 1, \end{cases}
\]
then the transition probabilities of the $y$-coordinate depend on the $x$-coordinate, so the projection is not Markov.

\begin{definition}
Given a Markov chain $\{X(n) : n \geq 0\}$ on $\mathcal{X}$, we say $\pi : \mathcal{X} \to \mathcal{Y}$ is a \textbf{Markov projection} if it is surjective and $\{\pi(X(n)) : n \geq 0\}$ is a Markov chain whose transition probabilities do not depend on the choice of initial conditions $X(0)$.
\end{definition}

\begin{remark}
Throughout these notes, all Markov chains are assumed time-homogeneous, following the convention of Kemeny and Snell.
\end{remark}

We address two questions:
\begin{enumerate}[label=\arabic*.]
\item What is a necessary and sufficient condition for $\pi$ to be a Markov projection?
\item If $\pi$ is a Markov projection, what are the transition probabilities of $Y(n) := \pi(X(n))$?
\end{enumerate}

\section*{Transition Probabilities of the Projected Process}

\begin{definition}[Notation]
Given a surjection $\pi : \mathcal{X} \to \mathcal{Y}$, let $V$ be the matrix with rows indexed by $\mathcal{X}$ and columns by $\mathcal{Y}$ defined by $V_{xy} = \mathbf{1}_{\{\pi(x) = y\}}$. Let $U$ be an arbitrary stochastic matrix with rows indexed by $\mathcal{Y}$ and columns by $\mathcal{X}$ satisfying $U_{yx} = 0$ whenever $\pi(x) \neq y$.
\end{definition}

Note that $V$ is stochastic and $U$ need not be unique. Neither matrix involves $X(n)$ or any probability.

\begin{theorem}
Suppose $\{X(n) : n \geq 0\}$ is a Markov chain with transition matrix $P$, and $\pi : \mathcal{X} \to \mathcal{Y}$ is a Markov projection. Then $\{Y(n) := \pi(X(n))\}$ has transition matrix
\[
\hat{P} = UPV.
\]
\end{theorem}

\begin{proof}
Intuitively: given $Y(n) = y'$, to find $P(Y(n+1) = y)$, since the probability cannot depend on $n$ or the initial conditions, one can arbitrarily choose $X(n) = x'$ (corresponding to $U$), let it update randomly to $X(n+1) = x$ (corresponding to $P$), then project to $y$ (corresponding to $V$).

For the formal proof, it suffices to consider $U$ with exactly one nonzero entry per row (the general case follows by linearity). Suppose $U_{y',x'} = 1$. Then
\[
(UPV)_{y'y} = \sum_{x \in \mathcal{X}} P_{x'x} V_{xy} = \sum_{x \in \pi^{-1}(y)} P_{x'x} = P(X(n+1) \in \pi^{-1}(y) \mid X(n) = x').
\]
Since $\pi$ is a Markov projection and $\pi(x') = y'$, this equals
\[
P(Y(n+1) = y \mid Y(n) = y'). 
\qedhere
\]
\end{proof}

\section*{Necessary and Sufficient Conditions for a Markov Projection}

\begin{theorem}
Suppose $\{X(n) : n \geq 0\}$ is a Markov chain on $\mathcal{X}$ and $\pi : \mathcal{X} \to \mathcal{Y}$ is surjective. Then $\pi$ is a Markov projection if and only if
\[
VUPV = PV
\]
for all $U$ and $V$ satisfying the conditions above.
\end{theorem}

\begin{remark}
Written as $V\hat{P} = PV$, this relation is called \textbf{Pitman-Rogers intertwining}.\footnote{L.C.G. Rogers, J.W. Pitman, ``Markov Functions,'' \textit{The Annals of Probability}, 9(4), 573--582, 1981.}
\end{remark}

\begin{proof}
($\Rightarrow$) Suppose $\pi$ is a Markov projection. On one hand,
\[
(PV)_{x',y} = \sum_{x \in \pi^{-1}(y)} P_{x'x} = P(X(1) \in \pi^{-1}(y) \mid X(0) = x').
\]
On the other hand,
\[
(VUPV)_{x',y} = \sum_{x \in \pi^{-1}(\pi(x'))} U_{\pi(x')x} \cdot P(X(1) \in \pi^{-1}(y) \mid X(0) = x).
\]
Since $\pi$ is a Markov projection, $P(X(1) \in \pi^{-1}(y) \mid X(0) = x) = P(X(1) \in \pi^{-1}(y) \mid X(0) = x')$ for all $x \in \pi^{-1}(\pi(x'))$. Factoring this out and using that the rows of $U$ sum to 1 gives $(VUPV)_{x',y} = (PV)_{x',y}$.

($\Leftarrow$) Suppose $VUPV = PV$ for all valid $U, V$. For any initial conditions $X(0)$, define $U$ by
\[
U_{yx} = \begin{cases} 0 & \text{if } \pi(x) \neq y \\ \frac{P(X(0)=x)}{P(X(0) \in \pi^{-1}(y))} & \text{if } \pi(x) = y \text{ and } P(X(0) \in \pi^{-1}(y)) \neq 0 \\ \mu(x) & \text{if } \pi(x) = y \text{ and } P(X(0) \in \pi^{-1}(y)) = 0 \end{cases}
\]
where $\mu(x)$ is an arbitrary probability measure on $\pi^{-1}(y)$. By assumption, $VUPV = PV$ does not depend on $X(0)$. Reusing the earlier calculation, this gives
\[
\sum_{x \in \pi^{-1}(\pi(x'))} \frac{P(X(0)=x)}{P(X(0) \in \pi^{-1}(x'))} P(X(1) \in \pi^{-1}(y) \mid X(0) = x).
\]
By the definition of conditional probability and countable additivity, this equals
\[
P(\pi(X(1)) = y \mid \pi(X(0)) = \pi(x')) = P(Y(1) = y \mid Y(0) = y'),
\]
which does not depend on the initial conditions $X(0)$. Therefore $\pi$ is a Markov projection.
\end{proof}

\begin{remark}
Rogers and Pitman generalize this to the case when $V$ is a stochastic matrix (which they call a Markov kernel).
\end{remark}

\end{document}